\documentclass{beamer}

% Theme Choice - "Madrid" is clean and professional for lectures
\usetheme{Madrid}
\usecolortheme{beaver}

% Packages
\usepackage{graphicx}
\usepackage{amsmath}
\usepackage{amssymb}
\usepackage{hyperref}

% Title Information
\title{Unit 3:National Income, Price Determination}
\subtitle{Aggregate Demand, Aggregate Supply, and Fiscal Policy}
\author{Doğukan Günkördü}
\date{\today}

\begin{document}

% -----------------------------------------------------------------------------
% Title Slide
% -----------------------------------------------------------------------------
\begin{frame}
    \titlepage
\end{frame}

% -----------------------------------------------------------------------------
% Table of Contents
% -----------------------------------------------------------------------------
\begin{frame}{Lecture Outline}
    \tableofcontents
\end{frame}

% -----------------------------------------------------------------------------
% Section 1: Aggregate Demand
% -----------------------------------------------------------------------------
\section{Aggregate Demand (AD)}

\begin{frame}{Aggregate Demand (AD)}
    \begin{itemize}
        \item \textbf{Definition}: The total quantity of output (Real GDP) demanded by all sectors of the economy (Households, Firms, Government, Foreigners) at various price levels.
        \item \textbf{The Curve}: Downward sloping.
        \item \textbf{Inverse Relationship}: As Price Level $\uparrow$, Real GDP Demanded $\downarrow$.
    \end{itemize}
    \vspace{0.5cm}
    \begin{block}{Why is AD Downward Sloping?}
        \begin{enumerate}
            \item \textbf{The Wealth Effect (Real Balances Effect)}: Higher prices reduce the purchasing power of accumulated assets, reducing consumption.
            \item \textbf{The Interest Rate Effect}: Higher prices increase the demand for money, driving up interest rates, which reduces Investment ($I$).
            \item \textbf{The Foreign Trade Effect}: Higher domestic prices make exports more expensive and imports cheaper, reducing Net Exports ($X_n$).
        \end{enumerate}
    \end{block}
\end{frame}

\begin{frame}{Shifters of Aggregate Demand}
    Any change in the components of GDP (excluding price level changes) shifts the AD curve.
    \[ AD = C + I + G + X_n \]
    \begin{itemize}
        \item \textbf{Consumption ($C$)}: Consumer confidence, taxes, wealth, debt.
        \item \textbf{Investment ($I$)}: Interest rates, business expectations, technology.
        \item \textbf{Government Spending ($G$)}: Changes in federal/state spending.
        \item \textbf{Net Exports ($X_n$)}: Foreign income, exchange rates.
    \end{itemize}
    \vspace{0.5cm}
    \textit{Key: Expansionary policy shifts AD Right ($\rightarrow$); Contractionary shifts Left ($\leftarrow$).}
\end{frame}

% -----------------------------------------------------------------------------
% Section 2: Multipliers
% -----------------------------------------------------------------------------
\section{Expenditure and Tax Multipliers}

\begin{frame}{Marginal Propensities}
    How much of an extra dollar do we spend vs. save?
    \begin{itemize}
        \item \textbf{Marginal Propensity to Consume (MPC)}: $\frac{\Delta \text{Consumption}}{\Delta \text{Income}}$
        \item \textbf{Marginal Propensity to Save (MPS)}: $\frac{\Delta \text{Savings}}{\Delta \text{Income}}$
    \end{itemize}
    \begin{block}{The Golden Rule}
        \[ MPC + MPS = 1 \]
    \end{block}
\end{frame}

\begin{frame}{The Spending and Tax Multipliers}
    \begin{block}{Spending Multiplier}
        Used for changes in Government Spending ($G$) or Investment ($I$).
        \[ \text{Multiplier} = \frac{1}{MPS} \quad \text{or} \quad \frac{1}{1 - MPC} \]
        \[ \Delta \text{Total AD} = \text{Multiplier} \times \Delta \text{Initial Spending} \]
    \end{block}

    \begin{block}{Tax Multiplier}
        Used for changes in lump-sum taxes.
        \[ \text{Tax Multiplier} = \frac{-MPC}{MPS} \]
        \textit{Note: The tax multiplier is always smaller than the spending multiplier because consumers save a portion of the tax cut.}
    \end{block}
\end{frame}

% -----------------------------------------------------------------------------
% Section 3: Short-Run Aggregate Supply
% -----------------------------------------------------------------------------
\section{Short-Run Aggregate Supply (SRAS)}

\begin{frame}{Short-Run Aggregate Supply (SRAS)}
    \begin{itemize}
        \item \textbf{Definition}: The relationship between the price level and quantity of goods supplied in the short run.
        \item \textbf{Slope}: Upward sloping.
    \end{itemize}
    \begin{block}{Why Upward Sloping? (Sticky Wages)}
        \begin{itemize}
            \item In the short run, input prices (wages, raw materials) are \textbf{sticky} (fixed by contracts).
            \item If Price Level rises, firms sell products for more money, but their costs stay the same.
            \item \textbf{Result}: Profit margins increase, so firms produce more.
        \end{itemize}
    \end{block}
\end{frame}

\begin{frame}{Shifters of SRAS}
    \textbf{Acronym: RAP} (Resources, Actions of Gov, Productivity)
    \begin{enumerate}
        \item \textbf{Resource Prices (Input Costs)}: 
        \begin{itemize}
            \item Wages, oil/energy prices, raw materials.
            \item Supply Shocks (negative or positive).
            \item Inflationary Expectations (if workers expect inflation, they demand higher wages $\rightarrow$ SRAS shifts left).
        \end{itemize}
        \item \textbf{Actions of Government}:
        \begin{itemize}
            \item Taxes on producers, subsidies, regulations.
        \end{itemize}
        \item \textbf{Productivity}:
        \begin{itemize}
            \item Technology, machinery, human capital.
        \end{itemize}
    \end{enumerate}
\end{frame}

% -----------------------------------------------------------------------------
% Section 4: Long-Run Aggregate Supply
% -----------------------------------------------------------------------------
\section{Long-Run Aggregate Supply (LRAS)}

\begin{frame}{Long-Run Aggregate Supply (LRAS)}
    \begin{itemize}
        \item \textbf{Definition}: The level of output the economy can produce when all resources are fully employed (Potential GDP / Full Employment Output $Y_f$).
        \item \textbf{Slope}: Vertical.
        \item \textbf{Rationale}: In the long run, all input prices (wages) are \textbf{flexible}. A change in the price level does not change real profit or output.
        \item \textbf{Shifters}:
        \begin{itemize}
            \item The same factors that shift the PPC (Production Possibilities Curve).
            \item Quantity/Quality of Land, Labor, Capital, and Technology.
        \end{itemize}
    \end{itemize}
\end{frame}

% -----------------------------------------------------------------------------
% Section 5: Equilibrium
% -----------------------------------------------------------------------------
\section{Equilibrium and Output Gaps}

\begin{frame}{Equilibrium in the AD-AS Model}
    \begin{itemize}
        \item \textbf{Short-Run Equilibrium}: Intersection of AD and SRAS. 
        \item \textbf{Long-Run Equilibrium}: Intersection of AD, SRAS, and LRAS.
    \end{itemize}
    \vspace{0.5cm}
    \textbf{Three States of the Economy:}
    \begin{enumerate}
        \item \textbf{Full Employment}: Equilibrium is on the LRAS line.
        \item \textbf{Recessionary Gap}: Equilibrium is to the \textbf{left} of LRAS. ($Y_e < Y_f$)
        \begin{itemize}
            \item High Unemployment.
        \end{itemize}
        \item \textbf{Inflationary Gap}: Equilibrium is to the \textbf{right} of LRAS. ($Y_e > Y_f$)
        \begin{itemize}
            \item Low Unemployment, High Inflation risk.
        \end{itemize}
    \end{enumerate}
\end{frame}

% -----------------------------------------------------------------------------
% Section 6: Self-Adjustment
% -----------------------------------------------------------------------------
\section{Long-Run Self Adjustment}

\begin{frame}{Long-Run Self Adjustment}
    If the government does nothing, how does the economy return to full employment?
    \vspace{0.3cm}
    \begin{block}{From a Recessionary Gap}
        \begin{enumerate}
            \item High unemployment exists.
            \item Nominal wages \textbf{fall} eventually (input costs decrease).
            \item SRAS shifts \textbf{Right}.
            \item Price Level falls, Output returns to $Y_f$.
        \end{enumerate}
    \end{block}

    \begin{block}{From an Inflationary Gap}
        \begin{enumerate}
            \item Labor market is tight (unemployment is too low).
            \item Nominal wages \textbf{rise} (workers demand more pay).
            \item SRAS shifts \textbf{Left}.
            \item Price Level rises, Output returns to $Y_f$.
        \end{enumerate}
    \end{block}
\end{frame}

% -----------------------------------------------------------------------------
% Section 7: Fiscal Policy (Detailed)
% -----------------------------------------------------------------------------
\section{Fiscal Policy}

\begin{frame}{Fiscal Policy Overview}
    \begin{itemize}
        \item \textbf{Definition}: The use of government spending ($G$) and taxation ($T$) to influence the economy's output, employment, and price level.
        \item \textbf{Discretionary Fiscal Policy}: Requires new legislation/action (e.g., Congress passing a stimulus bill).
        \item \textbf{Automatic Stabilizers}: Policies already in law that stabilize the economy without new action.
    \end{itemize}
\end{frame}

\begin{frame}{Expansionary Fiscal Policy}
    \textbf{Goal}: Combat \textbf{Recession} (Close a recessionary gap).
    \begin{block}{Tools}
        \begin{enumerate}
            \item \textbf{Increase Government Spending ($G \uparrow$)}: Directly increases AD.
            \item \textbf{Decrease Taxes ($T \downarrow$)}: Increases Disposable Income ($Y_d$), which increases Consumption ($C$).
        \end{enumerate}
    \end{block}
    \begin{itemize}
        \item \textbf{Mechanism}: $G \uparrow$ or $T \downarrow \Rightarrow AD \text{ shifts Right} \Rightarrow PL \uparrow, \text{Real GDP} \uparrow, \text{Unemployment} \downarrow$.
        \item \textbf{Side Effect}: Can lead to a budget deficit and higher price levels (inflation).
    \end{itemize}
\end{frame}

\begin{frame}{Contractionary Fiscal Policy}
    \textbf{Goal}: Combat \textbf{Inflation} (Close an inflationary gap).
    \begin{block}{Tools}
        \begin{enumerate}
            \item \textbf{Decrease Government Spending ($G \downarrow$)}: Directly decreases AD.
            \item \textbf{Increase Taxes ($T \uparrow$)}: Decreases Disposable Income ($Y_d$), which decreases Consumption ($C$).
        \end{enumerate}
    \end{block}
    \begin{itemize}
        \item \textbf{Mechanism}: $G \downarrow$ or $T \uparrow \Rightarrow AD \text{ shifts Left} \Rightarrow PL \downarrow, \text{Real GDP} \downarrow, \text{Unemployment} \uparrow$.
        \item \textbf{Side Effect}: Can lead to a budget surplus but may increase unemployment.
    \end{itemize}
\end{frame}

\begin{frame}{Spending vs. Tax Multiplier Nuances}
    Why is the Spending Multiplier effect stronger than the Tax Multiplier?
    \vspace{0.3cm}
    \begin{itemize}
        \item \textbf{Spending ($G$)}: Every dollar of government spending goes directly into the economy.
        \item \textbf{Taxes ($T$)}: A tax cut increases disposable income, but households \textbf{save} a portion of it (based on the MPS).
        \begin{itemize}
            \item Example: If $MPC = 0.8$, a $\$100$ tax cut leads to only $\$80$ in initial new spending ($20\%$ is saved).
        \end{itemize}
    \end{itemize}
    \vspace{0.3cm}
    \[ |\text{Spending Multiplier}| > |\text{Tax Multiplier}| \]
    \[ \frac{1}{MPS} > \frac{MPC}{MPS} \]
\end{frame}

\begin{frame}{Automatic Stabilizers}
    Mechanisms that dampen business cycle fluctuations \textit{without} new laws.
    
    \begin{block}{1. Progressive Income Tax System}
        \begin{itemize}
            \item \textbf{In a Recession}: Incomes fall $\rightarrow$ households move into lower tax brackets $\rightarrow$ tax burden falls automatically $\rightarrow$ cushions the drop in disposable income.
            \item \textbf{In an Expansion}: Incomes rise $\rightarrow$ households move into higher tax brackets $\rightarrow$ tax burden rises $\rightarrow$ slows down excessive consumption.
        \end{itemize}
    \end{block}

    \begin{block}{2. Transfer Payments (Unemployment Insurance)}
        \begin{itemize}
            \item \textbf{In a Recession}: More people apply for benefits $\rightarrow$ government spending ($G$) effectively increases $\rightarrow$ supports AD.
            \item \textbf{In an Expansion}: Fewer people qualify $\rightarrow$ spending decreases automatically.
        \end{itemize}
    \end{block}
\end{frame}

\begin{frame}{Issues with Fiscal Policy: Lags}
    Time delays can make fiscal policy less effective.
    \begin{enumerate}
        \item \textbf{Recognition Lag}: Time to realize there is a problem (collecting GDP data takes months).
        \item \textbf{Administrative Lag}: Time for Congress to pass the law.
        \item \textbf{Operational Lag}: Time for the money to actually be spent.
    \end{enumerate}
    \vspace{0.3cm}
    \textit{Risk: By the time policy hits, the economy may have already self-corrected.}
\end{frame}

\begin{frame}{Issues with Fiscal Policy: Crowding Out Effect}
    The negative impact of deficit spending on private investment.
    

[Image of Crowding Out Diagram]

    \begin{itemize}
        \item \textbf{Scenario}: Gov increases spending ($G \uparrow$) to fight recession, running a deficit.
        \item \textbf{The Chain Reaction}:
        \begin{enumerate}
            \item Government borrows money (sells bonds).
            \item Demand for Loanable Funds increases ($\uparrow$).
            \item Real Interest Rates ($r$) rise ($\uparrow$).
            \item Private Investment ($I$) decreases ($\downarrow$) because borrowing is expensive.
        \end{enumerate}
        \item \textbf{Result}: The AD shift is smaller than intended because $G \uparrow$ is partially offset by $I \downarrow$.
    \end{itemize}
\end{frame}

\begin{frame}{Unit 3 Summary}
    \begin{enumerate}
        \item \textbf{AD} is downward sloping (Wealth, Interest Rate, Trade effects).
        \item \textbf{SRAS} is upward sloping (Sticky Wages). \textbf{LRAS} is vertical (Flexible Wages).
        \item Equilibrium determines PL and Output.
        \item \textbf{Fiscal Policy} shifts AD active intervention.
        \item \textbf{Self-Adjustment} shifts SRAS passive correction.
        \item Watch out for \textbf{Crowding Out} when the government borrows!
    \end{enumerate}
\end{frame}

\end{document}
